% !TeX root = ../main.tex
\begin{frame}{课程评价——康鹤群}
  \begin{SlideCanvas}
    \node[anchor=north west,inner sep=0pt,text width=850pt,align=left,
      text=BICNavy] at (55,80) {\begin{minipage}{850pt}
      \usebeamerfont{normal text}\raggedright
      \setlength{\parindent}{0pt}\setlength{\parskip}{4pt}
      \clubpenalty=10000\widowpenalty=10000
\par\addvspace{6pt}{\usebeamerfont{BIC point}一、课程评价\par}\nobreak
课程涵盖核磁共振影像、人工听觉、医用机器人、柔性电子及超声换能器等主题，内容丰富，体现了医学与工程交叉融合的特点，有助于学生了解不同领域的发展方向和应用场景。\par
\par\addvspace{6pt}{\usebeamerfont{BIC point}二、学习收获\par}\nobreak
通过课程学习，我对医学影像、医疗器械和智能技术在临床中的应用有了更全面的认识，也更加理解医工交叉研究需要同时关注技术原理、实际需求和成果转化。\par
\par\addvspace{6pt}{\usebeamerfont{BIC point}三、意见与建议\par}\nobreak
希望今后的课程在介绍技术和应用的基础上，进一步增加学术深度。建议各专题围绕一个具体研究问题，讲解核心原理、研究方法与技术难点，并结合代表性论文或实验数据，分析现有成果的局限和改进方向。这样能帮助学生从了解前沿技术，进一步走向理解和评价相关研究，避免课程停留在普及性科普层面。\par
    \end{minipage}};
  \end{SlideCanvas}
  \note{康鹤群课程评价由用户提供，正文按原文保留，仅将标题与段落结构化。一、课程评价\par \par 课程涵盖核磁共振影像、人工听觉、医用机器人、柔性电子及超声换能器等主题，内容丰富，体现了医学与工程交叉融合的特点，有助于学生了解不同领域的发展方向和应用场景。\par \par 二、学习收获\par \par 通过课程学习，我对医学影像、医疗器械和智能技术在临床中的应用有了更全面的认识，也更加理解医工交叉研究需要同时关注技术原理、实际需求和成果转化。\par \par 三、意见与建议\par \par 希望今后的课程在介绍技术和应用的基础上，进一步增加学术深度。建议各专题围绕一个具体研究问题，讲解核心原理、研究方法与技术难点，并结合代表性论文或实验数据，分析现有成果的局限和改进方向。这样能帮助学生从了解前沿技术，进一步走向理解和评价相关研究，避免课程停留在普及性科普层面。}
\end{frame}
